\documentclass{article}
\usepackage[UTF8]{ctex}
\usepackage{amsmath,amssymb,amsfonts}
\usepackage{enumitem}
\usepackage{geometry}
\geometry{a4paper, margin=1in}
% === 错题本排版优化（不参与题目标识） ===
\setlength{\parindent}{0pt}
\setlength{\parskip}{0.35em plus 0.10em minus 0.08em}
\setlength{\emergencystretch}{3em}
\setlist{itemsep=0.25em, topsep=0.40em, parsep=0pt}
\allowdisplaybreaks[2]
% === 错题本排版优化结束 ===
\begin{document}
\section*{华中科技大学2025年高等代数真题}
\section*{124．华中科技大学2025年高等代数试卷}

\subsection{题目1}
1．（15 分）计算行列式

\[
D=\left|\begin{array}{ccccc}
1 & 1 & 1 & 1 & 1 \\
2021 & 2022 & 2023 & 2024 & 2025 \\
2021^{2} & 2022^{2} & 2023^{2} & 2024^{2} & 2025^{2} \\
2021^{3} & 2022^{3} & 2023^{3} & 2024^{3} & 2025^{3} \\
2021^{5} & 2022^{5} & 2023^{5} & 2024^{5} & 2025^{5}
\end{array}\right| .
\]

\subsection{题目2}
2．（15 分）设向量

\[
\alpha_{1}=(a, 2,10)^{\prime}, \alpha_{2}=(-2,1,5)^{\prime}, \alpha_{3}=(-1,1,4)^{\prime}, \beta=(1, b, 2)^{\prime} .
\]

当 \(a, b\) 满足什么条件时：

\begin{itemize}
  \item[（1）] \(\beta\) 不能表示为 \(\alpha_{1}, \alpha_{2}, \alpha_{3}\) 的线性组合？
  \item[（2）] \(\beta\) 可以唯一地表示为 \(\alpha_{1}, \alpha_{2}, \alpha_{3}\) 的线性表出？
  \item[（3）] \(\beta\) 可以以无穷种方式表示为 \(\alpha_{1}, \alpha_{2}, \alpha_{3}\) 的线性组合？写出所有的表示式。
\end{itemize}

\subsection{题目3}
3．（20 分）设 \(A, B, C, D\) 为 \(n\) 阶矩阵，且 \(C D^{\prime}=D C^{\prime}\) ，证明：

\[
\operatorname{det}\left(\begin{array}{ll}
A & B \\
C & D
\end{array}\right)=\operatorname{det}\left(A D^{\prime}-B C^{\prime}\right)
\]

\subsection{题目4}
4．（20 分）对于二阶矩阵 \(A, B_{1}, B_{2}, B_{3}, B_{4}, A \neq O\) ，若

\[
\operatorname{det}\left(A+B_{1}\right)=\operatorname{det}(A)+\operatorname{det}\left(B_{i}\right), i=1,2,3,4 .
\]

证明：矩阵 \(B_{1}, B_{2}, B_{3}, B_{4}\) 线性相关．

\subsection{题目5}
5．（20 分）设 \(A, B\) 为 \(n\) 阶实矩阵。证明：如果存在一个复可逆矩阵 \(P\) ，使得 \(P^{-1} A P=B\) 。那么一定存在一个实可逆矩阵 \(Q\) ，使得 \(Q^{-1} A Q=B\) 。

\subsection{题目6}
6．（20 分）设 \(A, B\) 为 \(n\) 阶矩阵，且 \(\operatorname{rank}(A)=\operatorname{rank}(B)\) 。证明：\(A^{2} B=A\) 的充要条件是 \(B^{2} A=B\) 。

\subsection{题目7}
7．（20 分）设 \(A\) 为实方阵，\(A+A^{\prime}\) 为正定矩阵，但 \(A \neq A^{\prime}\) 。证明： \(\operatorname{det}\left(A+A^{\prime}\right)<\operatorname{det}(2 A)\) 。

\subsection{题目8}
8．（20 分）设 \(\alpha, \beta, \gamma\) 是欧氏空间 \(\mathbb{R}^{n}\) 中三个非零向量，已知它们两两正交．记矩阵 \(A=\alpha \beta^{\prime}+\beta \gamma^{\prime}+\gamma \alpha^{\prime}\)．

\begin{itemize}
  \item[（1）] 证明： \(\operatorname{rank}(A)=3\) ．
  \item[（2）] \(A\) 是否可以相似对角化？请证明你的结论．
\end{itemize}
\end{document}
